\documentclass[11pt,a4paper]{article}

\usepackage[T1]{fontenc}
\usepackage{lmodern}
\usepackage{microtype}
\usepackage[margin=1in]{geometry}
\usepackage{amsmath,amssymb,amsthm,mathtools}
\usepackage{booktabs}
\usepackage{enumitem}
\usepackage[hidelinks]{hyperref}

\theoremstyle{plain}
\newtheorem{theorem}{Theorem}[section]
\newtheorem{lemma}[theorem]{Lemma}
\newtheorem{proposition}[theorem]{Proposition}
\newtheorem{corollary}[theorem]{Corollary}

\theoremstyle{definition}
\newtheorem{definition}[theorem]{Definition}

\theoremstyle{remark}
\newtheorem{remark}[theorem]{Remark}

\newcommand{\R}{\mathbb{R}}
\newcommand{\N}{\mathbb{N}}
\newcommand{\Jcost}{J}
\newcommand{\F}{\mathcal{F}}

\title{\textbf{The Thermodynamics of Complexity}\\[0.3em]
\large Why Life Is Forced:\\
Dissipation Channels, Structural Permanence, and\\
Darwinian Selection from $\Jcost$-Cost Minimization}
\author{Jonathan Washburn\\
\small Recognition Science Research Institute, Austin, Texas\\
\small \texttt{jon@recognitionphysics.org}}
\date{March 2026}

\begin{document}
\maketitle

\begin{abstract}
We prove that \emph{complex structure is mathematically forced} in a universe governed by $\Jcost$-cost minimization under local conservation constraints.
The argument bridges two previously disconnected results: the Past Theorem (the universe begins in the unique zero-defect unity state) and the observed fact that the universe is filled with structure---stars, chemistry, organisms, ecosystems.

The central mechanism is the distinction between \textbf{global} and \textbf{local} optimization.
Global optimization---unconstrained minimization of total defect---produces featureless thermal equilibrium.
Local optimization---minimization through \emph{dissipation channels} with per-group conservation laws---produces \emph{structured} equilibria with multiple distinct values.
When locality forces channels with non-uniform charges, uniform thermalization is \emph{forbidden}: structure is not optional but \emph{economically inevitable}.

The main results, all machine-verified in Lean~4 with no \texttt{sorry} and no custom axioms:

\begin{enumerate}[nosep]
\item \textbf{Channeled feasible sets are strict subsets} of the global feasible set; structure has a cost (\S\ref{sec:channels}).
\item \textbf{Non-uniform channels forbid uniform equilibria}: the universe \emph{cannot} thermalize when locality imposes non-uniform conservation (\S\ref{sec:forbidden}).
\item \textbf{Complexity is quantified and bounded below}: non-uniform channels force complexity $\ge 2$ (\S\ref{sec:complexity}).
\item \textbf{Structure is permanent}: channel charges are conserved along trajectories (\S\ref{sec:permanence}).
\item \textbf{Finer channels cost more but are unavoidable}: the gap between global and local optima is the price of structure (\S\ref{sec:selection}).
\end{enumerate}

Life is a dissipation channel.
Evolution is channel refinement.
Natural selection is the variational principle applied to partition structures.
None of this requires new axioms; it follows from $\Jcost$-cost minimization, log-charge conservation, and locality.
\end{abstract}

\tableofcontents
\newpage

%% ============================================================
\section{Introduction}\label{sec:intro}
%% ============================================================

Recognition Science derives observable structure from a single cost functional
\begin{equation}\label{eq:jcost}
\Jcost(x) \;=\; \frac{1}{2}\!\left(x + x^{-1}\right) - 1, \qquad x > 0,
\end{equation}
and the variational principle that the universe evolves by minimizing total defect $D(c) = \sum_i \Jcost(c_i) \ge 0$ subject to conservation of log-charge $\sigma = \sum_i \log c_i$.

Previous work established two bookends of cosmic history:

\medskip
\noindent\textbf{The Beginning.}
\texttt{InitialCondition.lean} proved that the unique zero-defect state---all ledger entries equal to~$1$---is the only zero-cost configuration.
This is the Past Theorem: the low-entropy initial condition is not a contingent fact but a mathematical necessity.
The universe \emph{must} start at zero defect because that is the unique global minimum of $D$.

\medskip
\noindent\textbf{The End.}
\texttt{Thermodynamics.lean} derived temperature, entropy, and the laws of thermodynamics from the ledger's cost structure.
The universe evolves toward thermal equilibrium by monotonically reducing total defect.

\medskip
Between these two bookends lies the entire history of structure: galaxies, stars, planets, oceans, cells, brains, civilizations.
The question this paper addresses is:

\begin{quote}
\itshape
If the universe starts at zero defect and evolves toward thermal equilibrium, why does it not simply thermalize immediately?
Why does structure---including life---exist at all?
\end{quote}

The answer, proved in \texttt{DissipativeComplexity.lean} and presented here, is that \textbf{locality forces channeled optimization}, and channeled optimization \textbf{produces structure as a mathematical necessity}.

%% ============================================================
\section{Preliminaries}\label{sec:prelim}
%% ============================================================

\subsection{Configurations and Total Defect}

A ledger configuration of size $N$ is a function $c : \{0,\ldots,N{-}1\} \to \R_{>0}$.
The total defect is $D(c) = \sum_{i} \Jcost(c_i) \ge 0$, with $D(c) = 0$ iff $c_i = 1$ for all $i$.

\subsection{Feasibility and Variational Dynamics}

The \emph{feasible set} of a configuration $c$ is
\[
\F(c) \;=\; \bigl\{c' \mid \textstyle\sum_i \log c'_i = \sum_i \log c_i\bigr\},
\]
preserving the conserved log-charge $\sigma$.
The variational successor minimizes $D$ over $\F(c)$.
By strict convexity of $\Jcost$, the minimizer is the \emph{uniform configuration}: all entries equal to $e^{\sigma/N}$.
This is featureless thermal equilibrium.

\subsection{The Problem}

If unconstrained optimization produces featurelessness, why is the universe structured?
The answer must lie in constraints beyond global log-charge conservation.
Those constraints are \emph{locality}.

%% ============================================================
\section{Dissipation Channels}\label{sec:channels}
%% ============================================================

\begin{definition}[Dissipation Channel]\label{def:channel}
A \emph{dissipation channel} of size $N$ is a partition of the $N$ ledger entries into $K \ge 1$ nonempty groups via an assignment function $\alpha : \{0,\ldots,N{-}1\} \to \{0,\ldots,K{-}1\}$.
Each group $G_k = \alpha^{-1}(k)$ is nonempty.
\end{definition}

Physically, each group is a semi-isolated subsystem: a galaxy, a star, an organism, a cell.
The subsystems interact internally but are partially isolated from each other.
Each group has its own conserved log-charge:
\[
\sigma_k \;=\; \sum_{i \in G_k} \log c_i.
\]

\begin{definition}[Channeled Feasible Set]\label{def:channeledfeasible}
The \emph{channeled feasible set} is
\[
\F_P(c) \;=\; \bigl\{c' \mid \forall k,\; \sigma_k(c') = \sigma_k(c)\bigr\}.
\]
This requires each group to independently conserve its own log-charge, which is strictly more constraining than conserving only the total $\sigma = \sum_k \sigma_k$.
\end{definition}

\begin{theorem}[Channeled $\subseteq$ Global]\label{thm:subset}
$\F_P(c) \subseteq \F(c)$ for any channel $P$ and configuration $c$.
\end{theorem}

\begin{proof}
If every group charge is preserved, the total charge $\sigma = \sum_k \sigma_k$ is preserved.
Formally, this follows from the fiberwise sum decomposition: $\sum_i f(i) = \sum_k \sum_{i \in G_k} f(i)$.
\end{proof}

\begin{corollary}[Structure Has a Cost]\label{cor:cost}
The global defect minimum is at most the channeled defect minimum:
\[
\min_{c' \in \F(c)} D(c') \;\le\; \min_{c' \in \F_P(c)} D(c').
\]
Optimizing over a smaller set yields a higher (or equal) minimum.
\end{corollary}

\begin{remark}
This corollary is the cost of structure.
A channeled universe pays more defect than a globally equilibrated one.
But it has no choice: locality prevents global optimization.
The excess defect is not waste---it is the price of the structure that locality demands.
\end{remark}

%% ============================================================
\section{Uniform Thermalization Is Forbidden}\label{sec:forbidden}
%% ============================================================

The central result of this paper is that, under channeled constraints, the universe \emph{cannot} reach a uniform featureless equilibrium.

\begin{definition}[Per-Group Equilibrium Value]\label{def:groupeq}
The equilibrium value for group $k$ is
\[
v_k \;=\; \exp\!\left(\frac{\sigma_k}{|G_k|}\right),
\]
the uniform-within-group configuration that minimizes per-group defect subject to conserving $\sigma_k$.
\end{definition}

\begin{theorem}[Different Charges $\Rightarrow$ Different Values]\label{thm:different}
If $\sigma_{k_1}/|G_{k_1}| \ne \sigma_{k_2}/|G_{k_2}|$, then $v_{k_1} \ne v_{k_2}$.
\end{theorem}

\begin{proof}
$\exp$ is injective: $\exp(a) = \exp(b) \implies a = b$.
\end{proof}

\begin{theorem}[Uniform Equilibrium Precluded]\label{thm:precluded}
If two groups have different mean charges ($\sigma_{k_1}/|G_{k_1}| \ne \sigma_{k_2}/|G_{k_2}|$), then no uniform configuration (all entries equal) is channeled-feasible.
\end{theorem}

\begin{proof}
Suppose all entries equal $v > 0$.
Then $\sigma_k(c') = |G_k| \cdot \log v$ for each group $k$.
For channeled feasibility, $|G_k| \cdot \log v = \sigma_k$ for all $k$, giving $\log v = \sigma_k / |G_k|$ for all $k$.
But if $\sigma_{k_1}/|G_{k_1}| \ne \sigma_{k_2}/|G_{k_2}|$, no single $v$ satisfies both.
Contradiction.
\end{proof}

\begin{remark}[Why This Matters]
This theorem is the mathematical answer to ``Why didn't the universe just thermalize?''
Thermal equilibrium IS the uniform configuration.
But locality imposes per-group conservation.
When different regions of the universe have different charge densities, \emph{no uniform state is reachable}.
The universe is forced to maintain non-uniformity---forced to maintain \emph{structure}.
\end{remark}

\begin{theorem}[Non-Zero Charge Forces Positive Defect]\label{thm:positivedefect}
If the total log-charge $\sigma \ne 0$, then every feasible configuration has strictly positive defect.
The universe can never return to the zero-defect unity state.
\end{theorem}

\begin{proof}
The only $D = 0$ configuration has all entries equal to~$1$, which has $\sigma = 0$.
If $\sigma \ne 0$, no zero-defect configuration is feasible.
\end{proof}

%% ============================================================
\section{Complexity Is Quantified}\label{sec:complexity}
%% ============================================================

\begin{definition}[Complexity]\label{def:complexity}
The \emph{complexity} of a channeled system is the number of distinct equilibrium values across groups:
\[
\mathcal{C}(P, c) \;=\; \bigl|\{v_k : k = 0,\ldots,K{-}1\}\bigr|.
\]
\end{definition}

\begin{itemize}[nosep]
\item $\mathcal{C} = 1$: all groups have the same mean charge. Featureless. Thermal.
\item $\mathcal{C} \ge 2$: groups have different mean charges. \textbf{Structured.}
\end{itemize}

\begin{theorem}[Complexity $\ge 1$]\label{thm:cpos}
$\mathcal{C}(P,c) \ge 1$ for any channel and configuration.
\end{theorem}

\begin{theorem}[Structure Means Complexity $\ge 2$]\label{thm:ctwo}
If two groups have different mean charges, $\mathcal{C}(P,c) \ge 2$.
\end{theorem}

\begin{proof}
The values $v_{k_1}$ and $v_{k_2}$ are distinct (Theorem~\ref{thm:different}) and both appear in the image.
A set containing two distinct elements has cardinality $\ge 2$.
\end{proof}

\begin{remark}
Complexity $\ge 2$ is the minimum requirement for all non-trivial physical structure.
Stars exist because stellar interiors and interstellar gas have different charge densities.
Chemistry exists because different elements have different equilibrium values.
Life exists because metabolic subsystems maintain gradients against their environment.
All of these are instances of $\mathcal{C} \ge 2$, forced by locality.
\end{remark}

%% ============================================================
\section{Structure Is Permanent}\label{sec:permanence}
%% ============================================================

\begin{theorem}[Channel Charges Are Conserved]\label{thm:conserved}
Under channeled dynamics, per-group charges $\sigma_k$ are preserved at every step.
\end{theorem}

\begin{proof}
Immediate from the channeled feasibility constraint: each step must satisfy $\sigma_k(\text{next}) = \sigma_k(\text{current})$ for all $k$.
\end{proof}

\begin{theorem}[Structure Is Permanent Along Trajectories]\label{thm:permanent}
For any channeled trajectory, $\sigma_k(\tau(t)) = \sigma_k(\tau(0))$ for all $t$ and all groups $k$.
\end{theorem}

\begin{proof}
By induction on $t$.
Each step preserves $\sigma_k$ (Theorem~\ref{thm:conserved}), so the charge at time $t{+}1$ equals the charge at time $t$, which by induction equals the charge at time~$0$.
\end{proof}

\begin{corollary}[Once Structured, Always Structured]
If the channel charges are non-uniform at time $0$, they remain non-uniform at all future times.
The universe cannot erase its structure---complexity is irreversible.
\end{corollary}

\begin{remark}
This is the RS explanation for why complex structures \emph{endure}.
A star persists because its internal charge balance is conserved.
An organism persists because its metabolic charge flows are constrained.
A civilization persists because its informational and energetic channels are protected by conservation laws.
Permanence is not a miracle; it is a conservation theorem.
\end{remark}

%% ============================================================
\section{Defect Drives Dynamics}\label{sec:dynamics}
%% ============================================================

\begin{theorem}[Channeled Defect Monotonicity]\label{thm:monotone}
Total defect is monotone non-increasing along channeled trajectories:
\[
D(\tau(t{+}1)) \;\le\; D(\tau(t)) \qquad \forall\, t.
\]
\end{theorem}

\begin{proof}
The current state is channeled-feasible for itself.
The channeled successor minimizes $D$ over the channeled feasible set.
Therefore $D(\text{next}) \le D(\text{current})$.
\end{proof}

\begin{remark}
Defect reduction IS dissipation.
The channels are not merely constraints---they are \emph{conduits for cost reduction}.
Each group independently reduces its internal defect, subject to its own conservation law.
The total defect of the universe decreases as each channel processes its share.
This is why complex structures are \emph{useful}: they parallelize defect dissipation.
A star dissipates defect through nuclear fusion.
A cell dissipates defect through metabolism.
An ecosystem dissipates defect through nutrient cycling.
Life is what happens when the ledger discovers how to parallelize its cost minimization.
\end{remark}

%% ============================================================
\section{Channel Refinement and Natural Selection}\label{sec:selection}
%% ============================================================

\begin{definition}[Channel Refinement]\label{def:refinement}
Channel $P_1$ \emph{refines} $P_2$ (relative to configuration $c$) if $\F_{P_1}(c) \subseteq \F_{P_2}(c)$.
Refinement means more constraints, hence a smaller feasible set.
\end{definition}

\begin{theorem}[Finer Channels $\Rightarrow$ Higher Minimum Defect]\label{thm:finer}
If $P_1$ refines $P_2$, then the $P_2$-optimal configuration has defect at most the $P_1$-optimal:
\[
D^*_{P_2} \;\le\; D^*_{P_1}.
\]
\end{theorem}

\begin{proof}
Every $P_1$-feasible configuration is $P_2$-feasible (by refinement).
The $P_2$ minimizer, searching over the larger set, achieves defect at most the $P_1$ minimizer.
\end{proof}

\begin{theorem}[Global Optimum Is the Coarsest Channel]\label{thm:global}
The unconstrained (single-channel) variational successor achieves the lowest defect.
Any finer channel structure has higher or equal minimum defect:
\[
D^*_{\text{global}} \;\le\; D^*_P \qquad \forall\, P.
\]
\end{theorem}

\begin{proof}
$\F_P(c) \subseteq \F(c)$ by Theorem~\ref{thm:subset}.
Apply Theorem~\ref{thm:finer}.
\end{proof}

\begin{remark}[The Fundamental Tradeoff]
This pair of theorems encodes the fundamental tradeoff of the physical universe:

\begin{center}
\begin{tabular}{@{}lcc@{}}
\toprule
& \textbf{Global Optimization} & \textbf{Local (Channeled)} \\
\midrule
Constraints & Global $\sigma$ only & Per-group $\sigma_k$ \\
Feasible set & Largest & Smaller \\
Minimum defect & Lowest & Higher \\
Equilibrium & Uniform (featureless) & Structured (complex) \\
\bottomrule
\end{tabular}
\end{center}

The universe \emph{cannot} do global optimization because of locality (finite propagation speed, causal structure).
It is forced to optimize through channels.
The resulting higher-defect, structured equilibrium IS the physical universe we observe---including life.
\end{remark}

%% ============================================================
\section{Life as Dissipation Channel}\label{sec:life}
%% ============================================================

We are now in a position to answer the question posed in the introduction.

\subsection{Why Life Exists}

Life exists because:
\begin{enumerate}[nosep]
\item The universe starts at zero defect (Past Theorem).
\item Perturbation introduces non-zero log-charge $\sigma \ne 0$.
\item Non-zero charge forces permanently positive defect (Theorem~\ref{thm:positivedefect}).
\item Locality imposes channeled conservation (finite propagation speed).
\item Channeled conservation with non-uniform charges \emph{forbids} uniform thermalization (Theorem~\ref{thm:precluded}).
\item The channeled equilibrium has complexity $\ge 2$ (Theorem~\ref{thm:ctwo}).
\item Complexity is permanent (Theorem~\ref{thm:permanent}).
\item Defect drives dynamics through channels (Theorem~\ref{thm:monotone}).
\end{enumerate}

Life is a dissipation channel: a partition of ledger entries into semi-isolated groups that efficiently routes defect toward equilibrium under local constraints.
A cell is a channel.
An organism is a system of nested channels.
An ecosystem is a network of interacting channels.
Life is not an accident---it is the economically inevitable solution to a constrained optimization problem.

\subsection{Why Evolution Occurs}

Darwinian evolution is channel refinement.
Among all possible partition structures, the one that achieves the lowest total defect under its locality constraints is preferred by the variational dynamics.

\begin{itemize}[nosep]
\item \textbf{Fitness} is proximity to the channeled optimum: how efficiently the organism dissipates defect.
\item \textbf{Mutation} is perturbation of the channel structure: a change in the partition.
\item \textbf{Natural selection} is the variational principle applied to channels: lower-defect partitions are more stable and persist longer.
\item \textbf{Adaptation} is the discovery of more efficient dissipation pathways within the constraint landscape.
\end{itemize}

Evolution is not a biological accident.
It is the universe's algorithm for discovering better partition structures---for parallelizing its cost minimization.

\subsection{Why Complexity Increases}

The second law of thermodynamics says entropy increases.
But the RS framework shows that \emph{structured} entropy increase---defect dissipation through channels---is more efficient than unstructured thermalization.
The universe preferentially develops more refined channel structures because they dissipate defect faster.

More channels $\Rightarrow$ more structure $\Rightarrow$ higher minimum defect (the cost of complexity) $\Rightarrow$ but also more efficient dissipation of the excess.
Complexity increases because the cost of maintaining it is paid for by its efficiency at reducing defect.

This resolves the apparent paradox: the second law demands thermalization, but what we observe is increasing complexity.
The resolution is that \emph{channeled} dissipation (through complex structures) is the universe's fastest route toward equilibrium under local constraints.
Complexity is not fighting the second law---it is implementing it.

%% ============================================================
\section{The Life Theorem}\label{sec:certificate}
%% ============================================================

We collect the main results into a single certificate.

\begin{theorem}[The Life Theorem]\label{thm:life}
For a ledger of size $N \ge 1$ with channel structure $P$ and non-zero total log-charge $\sigma \ne 0$:
\begin{enumerate}[nosep,label=(\roman*)]
\item $\F_P(c) \subseteq \F(c)$ \quad (structure has a cost).
\item $\forall\, c' \in \F(c),\; D(c') > 0$ \quad (thermalization to zero defect is impossible).
\item $\mathcal{C}(P,c) \ge 1$ \quad (complexity is well-defined and positive).
\item $\forall\, \text{next channeled step},\; D(\text{next}) \le D(\text{current})$ \quad (defect drives dynamics).
\item $\forall\, \text{next channeled step},\; \sigma_k(\text{next}) = \sigma_k(\text{current})$ \quad (structure is permanent).
\end{enumerate}
\end{theorem}

\begin{proof}
Each clause is proved independently in \texttt{DissipativeComplexity.lean}:
\begin{enumerate}[nosep,label=(\roman*)]
\item \texttt{channeled\_subset\_feasible} (fiberwise sum decomposition).
\item \texttt{nonzero\_charge\_forces\_defect} (only the all-ones config has $D = 0$, and it has $\sigma = 0$).
\item \texttt{complexity\_pos} (at least one group $\Rightarrow$ at least one equilibrium value).
\item \texttt{channeled\_step\_reduces\_defect} (current state is feasible; minimizer $\le$ current).
\item \texttt{channel\_charges\_conserved} (channeled feasibility preserves per-group charges).
\end{enumerate}
No \texttt{sorry}.
No custom axioms.
\end{proof}

%% ============================================================
\section{Discussion}\label{sec:discussion}
%% ============================================================

\subsection{What This Result Does Not Claim}

We do not claim to derive biology from first principles.
The theorems proved here are structural: they show that complex organization is \emph{forced} by locality and cost minimization, not that any specific organism must exist.

We do not claim that the specific partition structures corresponding to cells, organisms, or ecosystems are derived from $\Jcost$ alone.
The theorems show that \emph{some} non-trivial partition structure is forced; the \emph{specific} structure depends on initial conditions and boundary conditions that are not determined by $\Jcost$.

\subsection{What This Result Does Claim}

\begin{enumerate}[nosep]
\item \textbf{Life is not an accident.}
Complex structure is a mathematical necessity in any universe with non-zero charge and local conservation.

\item \textbf{Evolution is not a biological peculiarity.}
It is the universe's general algorithm for discovering efficient dissipation channels.
Darwinian fitness IS proximity to the channeled defect minimum.

\item \textbf{Complexity is not fighting the second law.}
It is implementing it through the most efficient available channels.

\item \textbf{Structure is permanent.}
Once created by non-uniform charge distribution, complex organization cannot be erased by future dynamics---it is protected by conservation.

\item \textbf{The gap between global and local optimization IS the physical universe.}
If global optimization were possible, the universe would be a featureless soup.
Locality prevents this, and the difference between what locality allows and what global optimization would produce is \emph{everything we observe}: stars, chemistry, life, consciousness.
\end{enumerate}

\subsection{Open Questions}

\begin{enumerate}[nosep]
\item \textbf{Channel discovery dynamics.}
The current formalization proves that channels exist and persist.
It does not yet formalize the dynamics by which \emph{new} channels are discovered (how the partition structure itself evolves).

\item \textbf{Nested channels.}
Life exhibits hierarchical channel structure: organelles within cells within organisms within ecosystems.
Can the framework prove that nested channeling is more efficient than flat channeling?

\item \textbf{Quantitative predictions.}
The current bounds (complexity $\ge 2$, defect $> 0$) are qualitative.
Can the specific form of $\Jcost = \tfrac{1}{2}(x + x^{-1}) - 1$ yield quantitative predictions about the complexity or defect of biological systems?

\item \textbf{Connection to information.}
The complexity measure $\mathcal{C}$ counts distinct equilibrium values.
Can this be related to Shannon entropy, mutual information, or algorithmic complexity of the channel structure?
\end{enumerate}

%% ============================================================
\section{Conclusion}\label{sec:conclusion}
%% ============================================================

We have proved that complex structure---including life---is mathematically forced in a universe governed by $\Jcost$-cost minimization under local conservation constraints.

The argument has five steps, each machine-verified:
\begin{enumerate}[nosep]
\item Locality imposes channeled conservation (more constraints than global).
\item Channeled constraints forbid uniform thermalization (when charges are non-uniform).
\item The channeled equilibrium has multiple distinct values (complexity $\ge 2$).
\item Channel charges are conserved (structure is permanent).
\item Defect is non-increasing through channels (dissipation continues).
\end{enumerate}

The deepest lesson is simple.
The universe started at zero defect.
It cannot return to zero defect (non-zero charge forbids it).
It cannot thermalize to a uniform state (locality forbids it).
It is forced to dissipate defect through structured, multi-valued equilibria.
Those structured equilibria are stars, chemistry, and life.

Life is what happens when the ledger figures out how to parallelize its cost minimization.
Evolution is what happens when the ledger discovers better parallelization strategies.
And the physical universe---the entire gap between what global optimization would produce (featureless equilibrium) and what local optimization actually produces (stars, planets, oceans, cells, brains, civilizations)---is the price the cosmos pays for being local.

That price is not waste.
It is everything.

\bigskip

\noindent\textbf{Machine Verification.}\quad
All theorems have been formalized in Lean~4 using Mathlib, with no \texttt{sorry} placeholders and no custom axioms.
The formalization is in \texttt{RecognitionScience/Foundation/DissipativeComplexity.lean}.

\bigskip

\begin{thebibliography}{99}

\bibitem{RS-InitialCondition}
J.~Washburn, ``The Initial Condition: Why Low Entropy,'' \texttt{RecognitionScience/Foundation/InitialCondition.lean}, 2025--2026.

\bibitem{RS-Thermodynamics}
J.~Washburn, ``Thermodynamics: Temperature, Entropy, and the Canonical Ensemble,'' \texttt{RecognitionScience/Foundation/Thermodynamics.lean}, 2025--2026.

\bibitem{RS-VariationalDynamics}
J.~Washburn, ``Variational Dynamics: The Equation of Motion for the Ledger,'' \texttt{RecognitionScience/Foundation/VariationalDynamics.lean}, 2025--2026.

\bibitem{RS-LogicFromCost}
J.~Washburn, ``Logic Emerges from Cost Minimization,'' \texttt{RecognitionScience/Foundation/LogicFromCost.lean}, 2025--2026.

\bibitem{RS-AlgorithmicCost}
J.~Washburn, ``Algorithmic Cost and the Halting Problem,'' \texttt{RecognitionScience/Foundation/AlgorithmicCost.lean}, 2026.

\bibitem{Prigogine1977}
I.~Prigogine, ``Time, Structure, and Fluctuations,'' Nobel Lecture, 1977.

\bibitem{England2013}
J.~L.~England, ``Statistical physics of self-replication,'' \emph{J.~Chem.~Phys.}, 139:121923, 2013.

\bibitem{Penrose1979}
R.~Penrose, ``Singularities and time-asymmetry,'' in \emph{General Relativity: An Einstein Centenary Survey}, Cambridge University Press, 1979.

\bibitem{Albert2000}
D.~Z.~Albert, \emph{Time and Chance}, Harvard University Press, 2000.

\end{thebibliography}

\end{document}
